%=============================================================================
% CHAPTER 6: CONCLUSION
%=============================================================================

\clearpage

\chapter{Conclusion}

\section{Summary of Achievements}

The aim of this thesis was to design, implement, and evaluate a hybrid password strength evaluation tool. This aim has been achieved. The tool combines five complementary analysis dimensions---rule-based length scoring, character diversity assessment, Shannon entropy estimation, pattern detection (including reversed keyboard patterns), and dictionary comparison with Levenshtein similarity---into a single transparent 0--100 scoring model with qualitative labels and actionable feedback.

The implementation is a Python-based desktop application with a CustomTkinter graphical user interface that provides real-time password analysis through a 300\,ms debounced input mechanism. The two-column layout presents a semicircular gauge and numerical score alongside detailed analysis results and improvement suggestions. A cryptographically secure password generator using Python's \texttt{secrets} module is integrated, allowing users to create strong passwords with a single click.

The tool was developed through an iterative process of 18 systematic enhancements and a user-experience redesign, resulting in a polished application that balances analytical rigour with ease of use. The modular architecture ensures that the evaluation logic is independent of the graphical interface and can be tested, modified, or extended in isolation.

\clearpage

\section{Final Conclusions}

Based on the development experience, functional testing, and comparative analysis, I draw the following conclusions:

The hybrid approach provides a reasonable balance of accuracy and simplicity for an educational tool. No single evaluation method is sufficient on its own---rule-based checks miss patterns, entropy misses dictionary words, and dictionary checks miss structural weaknesses. By combining all three, the tool provides a more comprehensive assessment than any single method alone.

Real-time feedback is superior to manual-triggered analysis. Users see results as they type, which encourages iterative improvement of their passwords. The 300\,ms debounce is fast enough to feel instantaneous while avoiding unnecessary computation during rapid typing. This design eliminates the need for a separate ``Analyze'' button, reducing the number of steps required to evaluate a password.

Named scoring constants make the model transparent and auditable, addressing a key shortcoming of ``black box'' strength meters. Users and developers can understand exactly how the score is computed by examining the named constants in the evaluator module, and the scoring weights can be adjusted by changing only the constant definitions.

The iterative improvement process (18 enhancements + UX redesign) demonstrates that even small tools benefit from systematic quality refinement. The most impactful improvements were the addition of the Levenshtein similarity check (which catches leetspeak variants), the automatic generation of reversed keyboard patterns (which catches patterns like \texttt{ytrewq}), and the user-experience redesign to the two-column layout (which eliminated the need for scrolling and improved the information density).

Reverse keyboard pattern detection is a novel contribution that most existing online checkers lack. While the individual value of this feature is modest (few users deliberately type \texttt{ytrewq}), it demonstrates that a hybrid approach can identify weaknesses that purely entropy-based or rule-based checkers miss. More broadly, it illustrates the principle that pattern detection should consider not only the most common forms of a pattern but also their transformations.

The tool's distinctive value lies in the combination of three principles: \textit{transparency}, \textit{privacy}, and \textit{educational feedback}. The per-dimension score breakdown makes the evaluation process fully transparent---users can see exactly how each dimension contributes to the final score, and the toggleable design ensures that this information is available on demand without cluttering the interface for casual users. The local-only, offline processing model provides a strong privacy guarantee that no online checker can match, making the tool suitable for evaluating even the most sensitive passwords. The contextual, natural-language suggestions---generated within the evaluation engine to ensure a single source of truth---educate users about the reasoning behind each weakness, promoting lasting improvements in password hygiene rather than one-time compliance.

Taken together, these principles make the tool not merely a password checker but an educational instrument that empowers users to make informed decisions about their digital security. In a landscape where most password tools provide only a binary ``weak/strong'' rating, this tool offers understanding.

\clearpage

\section{Future Work}

Several directions for future development have been identified:

\begin{itemize}
    \item \textbf{Online breach database integration:} Incorporating the Have I Been Pwned API (or a local k-anonymity-based query) would allow the tool to check passwords against billions of compromised credentials in real time while preserving privacy. This would dramatically improve the dictionary component's coverage.

    \item \textbf{Machine-learning-based scoring:} Training a neural network on large password datasets (e.g., the RockYou leak) could provide more nuanced strength predictions that account for frequency-based character distributions and common password structures, as demonstrated by Melcher et al.\ \cite{melcher2016}.

    \item \textbf{Multi-language dictionaries:} Adding common-password lists in multiple languages (Spanish, French, German, etc.) would improve the dictionary component's coverage for non-English-speaking users. Culture-specific patterns (e.g., keyboard layouts for different languages such as AZERTY and QWERTZ) could also be incorporated.

    \item \textbf{Passphrase support:} Adding specific evaluation logic for multi-word passphrases, including detection of well-known phrases and quotations, would extend the tool's applicability to the growing adoption of passphrase-based authentication.

    \item \textbf{Systematic leetspeak normalisation:} Implementing a substitution table that normalises common leetspeak substitutions (e.g., @ $\to$ a, 3 $\to$ e, 0 $\to$ o, \$ $\to$ s) before dictionary comparison would catch more sophisticated variants that the current Levenshtein approach misses.

    \item \textbf{Web and mobile deployment:} Wrapping the evaluation engine in a Flask or Django web application, or porting the core logic to JavaScript for browser-based deployment, would make the tool accessible to a wider audience. A React Native port could bring the tool to mobile platforms.

    \item \textbf{Configurable scoring weights:} Adding UI controls (e.g., sliders) that allow users to adjust the relative weights of the five scoring dimensions would make the tool useful for educational exploration, helping students understand how different factors contribute to password strength.

    \item \textbf{Larger common-password list:} Expanding the dictionary from 181 entries to 10,000 or more entries (sourced from publicly available breach datasets) would significantly improve the dictionary component's detection rate.

    \item \textbf{Responsive layout:} Making the window resizable and the layout responsive to different screen sizes would improve usability on laptops with smaller displays.

    \item \textbf{Integration into security training:} Packaging the tool as a component of organisational security training platforms could help raise password awareness in corporate and educational settings.
\end{itemize}